\section{Descrizione del prodotto}

\subsection{Obiettivi}
Il progetto \textbf{Sonar} nasce dall'esigenza dell'azienda proponente \textit{Halue S.r.l. Società Benefit} di realizzare una piattaforma web per la generazione automatica e ripetibile della documentazione vivente di store \textit{Shopify} di dimensioni \textit{enterprise}. \\
Nel corso del loro ciclo di vita, gli store di grandi dimensioni tendono a trasformarsi in sistemi stratificati, accumulando applicazioni di terze parti, \textit{webhook}, definizioni di \textit{metafield}, \textit{metaobject} e frammenti di codice di tema di cui si perde progressivamente traccia. \\
\\
Il principio cardine del sistema è che l'unica documentazione affidabile nel tempo non può dipendere dalla disciplina manuale degli sviluppatori, bensì deve essere ricostruita automaticamente dal sistema stesso a ogni scansione. L'obiettivo di \textbf{Sonar} non si limita a fornire un inventario statico di ciò che esiste nello store, ma mira a ricostruire l'intera rete delle dipendenze reciproche per rispondere a quesiti architetturali critici (es. valutare l'impatto della rimozione di una risorsa) e identificare componenti orfani o codice morto.

\subsection{Funzionalità}
L'architettura funzionale di \textbf{Sonar} si articola in cinque fasi sequenziali, progettate per mantenere una netta separazione delle responsabilità (\textit{separation of concerns}):
\begin{enumerate}
    \item \textbf{Acquisizione (Scansione)}\\
    Esecuzione di interrogazioni durevoli e resilienti verso le API di Shopify (GraphQL e REST) e lettura dei file del tema, finalizzate a produrre uno \textit{snapshot}, ovvero una fotografia immutabile e datata dello stato dello store. \\
    Per superare i limiti strutturali delle API di Shopify su store di grandi dimensioni, il processo gestisce nativamente:
    \begin{itemize}
        \item \textbf{Paginazione tramite cursori:} per l'estrazione incrementale di elenchi molto estesi senza superare i limiti di \textit{rate limiting};
        \item \textbf{Triangolazione delle fonti:} poiché alcune entità non sono visibili da un unico endpoint (ad esempio, webhook registrati da app terze), il sistema incrocia fonti eterogenee per massimizzare la copertura dell'inventario;
        \item \textbf{Telemetria interna:} registrazione puntuale delle risorse consumate per ogni singola scansione (chiamate API, passi di workflow, token LLM).
    \end{itemize}
    \item \textbf{Normalizzazione e modello a grafo} \\
    Trasformazione dei dati grezzi raccolti in un modello a grafo unificato. I \textbf{nodi} rappresentano le entità strutturali dello store (applicazioni, webhook, definizioni di metafield, file di tema, funzioni, canali di vendita), mentre gli \textbf{archi} codificano le dipendenze dirette e indirette tra di essi. \\
    Ogni arco è corredato dalla fonte dell'evidenza e da un \textbf{livello di confidenza} esplicito (\textit{accertato}, \textit{probabile} o \textit{ipotizzato});
    \item \textbf{Analisi statica del codice del tema}\\
    Ispezione isolata e sicura dei file del tema (in particolare template, sezioni e snippet scritti in linguaggio \textit{Liquid}) volta a rilevare i riferimenti ai metafield e a ricostruire il grafo di inclusione tra i file;
    \item \textbf{Rilevamento degli orfani e del codice morto}\\
    Applicazione di un motore di regole estensibile sul grafo normalizzato e sui risultati dell'analisi statica per individuare entità scollegate, definizioni di metafield mai referenziate né popolate, file di tema non inclusi e configurazioni obsolete;
    \item \textbf{Visualizzazione ed esportazione}\\
    Presentazione dei risultati tramite un'interfaccia web interattiva che consente di esplorare l'inventario, navigare visivamente il grafo delle dipendenze, filtrare le anomalie per livello di confidenza ed esportare la documentazione generata.
\end{enumerate}

\subsubsection{Vincoli e limiti di dominio}
Al fine di delimitare con precisione i confini operativi del sistema ed evitare ambiguità progettuali, sono stabiliti i seguenti vincoli di esclusione:
\begin{itemize}
    \item \textbf{Operatività strettamente in sola lettura}\\
    Sonar non apporta alcuna modifica agli store analizzati: non disinstalla applicazioni, non rimuove metafield, non altera il codice del tema. Il sistema si limita a segnalare le anomalie e a produrre istruzioni operative per eventuali interventi manuali da parte dei tecnici, richiedendo esclusivamente permessi (\textit{scope}) di sola lettura;
    \item \textbf{Assenza di monitoraggio in tempo reale}\\
    L'acquisizione dello stato dello store avviene unicamente tramite scansioni avviate manualmente su richiesta o secondo pianificazione periodica, senza flussi continui in tempo reale;
    \item \textbf{Esclusione dell'analisi prestazionale}\\
    Il sistema non rileva né valuta metriche di performance dello store o dello storefront;
    \item \textbf{Nessuna gestione dei dati di business}\\
    Entità commerciali quali prodotti, ordini e anagrafiche clienti non rientrano nel dominio applicativo di Sonar. Tali risorse vengono consultate unicamente in forma aggregata e statistica come evidenza per verificare l'effettivo utilizzo dei metafield configurati.
\end{itemize}

\subsection{Caratteristiche dell'utente}
Il prodotto si rivolge a figure professionali con competenze tecniche che operano nella gestione, manutenzione ed evoluzione di ecosistemi e-commerce complessi, quali:
\begin{itemize}
    \item \textbf{Sviluppatori Shopify e Full-Stack:} incaricati di intervenire sul codice del tema, integrare nuove applicazioni o effettuare attività di refactoring e pulizia del debito tecnico;
    \item \textbf{Architetti Software e Tech Lead:} responsabili della supervisione architetturale dello store, dell'audit tecnico e della valutazione dei rischi prima di rilasci evolutivi o migrazioni;
    \item \textbf{Project Manager tecnici e Solution Architect:} interessati a ottenere rapidamente una mappatura oggettiva e documentata dello stato di salute di uno store in fase di presa in carico (\textit{onboarding}) o consulenza.
\end{itemize}
Tali utenti possiedono familiarità con le architetture web e con i concetti base di Shopify, ma necessitano di uno strumento automatizzato che riduca i tempi di ispezione manuale e prevenga errori derivanti da dipendenze nascoste.

\subsection{Tecnologie}
In conformità con i vincoli architetturali imposti dal capitolato d'appalto, l'ecosistema tecnologico del prodotto si fonda sulle seguenti componenti:
\begin{itemize}
    \item \textbf{Framework applicativo:} \textbf{Next.js} per lo sviluppo dell'interfaccia web e delle API interne;
    \item \textbf{Infrastruttura di esecuzione e orchestrazione:} ecosistema \textbf{Vercel}, articolato in:
    \begin{itemize}
        \item \textit{Vercel Workflows} e \textit{Queues} per garantire l'esecuzione durevole delle scansioni e la ripresa automatica in caso di rate limiting o interruzioni;
        \item \textit{Vercel Sandbox} per l'esecuzione isolata e sicura dell'analisi statica sui file del tema;
        \item \textit{Vercel Blob} (per l'archiviazione degli snapshot grezzi).
    \end{itemize}
    \item \textbf{Persistenza e sicurezza dei dati:} \textbf{Supabase} (PostgreSQL) per la memorizzazione relazionale delle entità, del grafo delle dipendenze e delle metriche di scansione, protetto tramite politiche di \textit{Row Level Security (RLS)} per garantire l'isolamento multi-tenant dei dati;
    \item \textbf{Autenticazione verso Shopify:} gestione sicura delle credenziali di connessione agli store, inclusa l'archiviazione cifrata e il rinnovo automatico dei token di accesso generati tramite flusso \textit{client credentials}, soggetti a scadenza ogni 24 ore secondo le policy Shopify introdotte dal 2026.
\end{itemize}
